% =====================================================================
%  00_resumen.tex — Resumen del documento (va antes del índice)
% =====================================================================
\chapter*{Resumen}
\addcontentsline{toc}{chapter}{Resumen}

El control de entradas y salidas de las unidades del Metrobús en el Patio
CCA depende del registro manual, lo que provoca movimientos perdidos, horas
poco confiables y falta de evidencia cuando un registro se cuestiona. El
proyecto \saru{} (\emph{Sistema Automatizado de Reconocimiento de
Unidades}), desarrollado por Grupo CISA, propone un sistema de visión
artificial que detecta cada unidad en la puerta del patio, lee su número
económico, decide si entra o sale mediante una línea virtual, registra el
movimiento y avisa a los responsables. Todo el procesamiento ocurre en el
borde, en una computadora NVIDIA Jetson Orin, con tres modelos YOLO en
cascada optimizados con TensorRT.

El desarrollo sigue la metodología de \textbf{prototipado evolutivo}. Este
documento reporta el \textbf{Prototipo~1}: la cadena completa de
detección, lectura, cruce, registro, avisos, tablero web y grabación
continua, que opera en vivo en el patio desde noviembre de 2026. Se
presentan sus requerimientos, historias de usuario, arquitectura, protocolo
de pruebas y el trabajo previsto para los siguientes prototipos.
